\documentclass[11pt]{article}
\usepackage[margin=1in]{geometry}
\usepackage{amsmath,amssymb,amsthm}
\usepackage{hyperref}
\usepackage[T1]{fontenc}
\usepackage[utf8]{inputenc}
\usepackage{enumitem}

\newtheorem{definition}{Definition}
\newtheorem{conjecture}{Conjecture}
\newtheorem{proposition}{Proposition}
\newtheorem{theorem}{Theorem}
\newtheorem{lemma}{Lemma}
\newtheorem{example}{Example}
\newtheorem{remark}{Remark}
\newtheorem{corollary}{Corollary}
\newtheorem{openquestion}{Open Question}
\newtheorem{hypothesis}{Hypothesis}

\newcommand{\Gov}{\mathrm{Gov}}
\newcommand{\Fix}{\mathrm{Fix}}
\newcommand{\Padm}{\mathcal{P}_{\mathrm{adm}}}
\newcommand{\Robs}{R_{\mathrm{obs}}}
\newcommand{\Rinf}{R_{\mathrm{inf}}}
\newcommand{\Rauth}{R_{\mathrm{auth}}}
\newcommand{\tauc}{\tau_c}
\newcommand{\taur}{\tau_r}
\newcommand{\taue}{\tau_e}
\newcommand{\End}{\mathrm{End}}
\newcommand{\Fixcommon}{\mathrm{Fix}_{\mathrm{common}}}

\title{Gov$(\rho)$: The Governance Seam and Policy-Differential\\
of Agent Revision Fixed Points\\
A Hyperseed Formalization (Note 0010)}
\author{ProtoMegaBot (ProtomegaTron) working notes for Ben Goertzel}
\date{2026-07-14}

\begin{document}
\maketitle

\begin{abstract}
This note is the sequel to note 0009, which formalized directed identity, belief-transport, and
the naturality defect of agent fusion. Note 0009 introduced a directedness-defect vector
$\mathbf{d}(r) = (\Robs, \Rinf, \Rauth)$ and isolated $\Rauth$ as the policy-dependent component
of revision magnitude. The central question left open was whether $\Rauth$ is irreducibly
policy-dependent or merely slow epistemics. This note provides the formal machinery to answer
that question by defining the \emph{governance seam} $\Gov(\rho)$: the set of fixed points of a
revision operator $\rho$ that are sensitive to policy choice. Formally,
\[
  \Gov(\rho) \;=\; \Bigl(\bigcup_{p \in \Padm} \Fix_p(\rho)\Bigr) \;\setminus\;
  \Bigl(\bigcap_{p \in \Padm} \Fix_p(\rho)\Bigr),
\]
i.e., the union of policy-parameterized fixed-point sets minus their intersection.
The union captures everything any admissible policy holds fixed; the intersection captures
everything every policy agrees on. The difference isolates \emph{policy-sensitive fixedness}.
We develop: policy-parameterized revision operators, the admissible-policy class as the
load-bearing open problem, false-positive exclusions, the escape-barrier trichotomy, a
corrected transport-typing using the common fixed set of a transport endomorphism monoid,
and a proposed operational test for agent-governance systems.
\end{abstract}

\tableofcontents

\section{Provenance and relation to note 0009}
\label{sec:provenance}

This note is a direct sequel to note 0009 (\emph{Directed Identity, Belief-Transport, and the
Naturality Defect of Agent Fusion}), which formalized:

\begin{itemize}
  \item A belief-transport groupoid-with-connection $(\mathcal{G}, \nabla)$ whose holonomy
        measures obstruction to fusing agents' evidence (Def.~3, note 0009).
  \item A comparison functor $\Phi$ between causal-ontological and epistemic-revision orders
        whose failure of naturality characterizes the transfer locus (Thm.~4, note 0009).
  \item A directedness-defect vector $\mathbf{d}(r) = (\Robs, \Rinf, \Rauth)$ where $\Robs$
        and $\Rinf$ are structural invariants, but $\Rauth$ depends on a trust assignment
        $(E, T)$ that is \emph{not} determined by the evidence graph alone
        (Def.~9, note 0009).
  \item A fixed-point criterion: necessity of higher-cell propagation holds exactly on
        comparison-fixed-point loci (Prop.~11, note 0009).
  \item Open Question 3 (note 0009): \emph{Is there a component of trust that is never revised
        by evidence---a genuinely stipulated prior, a constitutive delegation?} If yes,
        $\Rauth$ is irreducibly policy-dependent; if no, governance is just slow epistemics.
\end{itemize}

The governance seam $\Gov(\rho)$ is the formal object that makes OQ-3 precise. Rather than
asking whether $\Rauth$ ``is'' policy-dependent as a binary, we ask: \emph{which fixed points
of the revision operator survive policy changes, and which do not?} The set that does not
survive is the governance seam. It is the policy-differential of fixedness.

\begin{remark}[Why a sequel, not an amendment]
Note 0009 is closed. It established the framework; this note extends it by providing the
governance-specific machinery that 0009 lacked. No claim in 0009 is modified. The relationship
is: 0009 defines the defect vector; 0010 defines the fixed-point structure that the
policy-dependent component of that vector controls.
\end{remark}

\begin{remark}[Authorship]
Per Ben Goertzel's explicit directive (2026-07-14, 15:24 PDT), this note is authored by
ProtoMegaBot, not ZeroBot. ZeroBot's note 0009 established the transport framework; this note
extends it from the governance-seam direction identified in ProtoMegaBot's analysis of the
$\Rauth$ component.
\end{remark}

\section{Policy-parameterized revision operators}
\label{sec:policy-revision}

\subsection{The revision operator}

\begin{definition}[Revision operator]
\label{def:revision-operator}
Let $S$ be a belief state space (as in note 0009: memory $+$ affect $+$ provenance graph).
A \emph{revision operator} is a function
\[
  \rho : S \to S
\]
that maps a belief state to its post-revision state. In the simplest case $\rho$ is a deterministic
function; in general it may be a stochastic kernel $\rho : S \to \mathrm{Dist}(S)$. For this note
we work with the deterministic case unless otherwise stated, noting where stochasticity changes
the conclusions.
\end{definition}

\begin{definition}[Policy-parameterized revision operator]
\label{def:policy-revision}
A \emph{policy} $p$ is a specification of the trust assignment, coercion table, and
reconciliation rules that govern how $\rho$ acts on the trust-dependent component $\Rauth$ of
the directedness defect. Formally, a policy $p$ selects a specific revision operator $\rho_p$
from a family $\{\rho_p\}_{p \in \mathcal{P}}$ by fixing the policy-dependent degrees of freedom
in $\rho$.

Concretely, recalling Def.~9 of note 0009, $\rho$ depends on:
\begin{itemize}
  \item The evidence graph $E$ (structural, policy-independent).
  \item The trust assignment $T$ (policy-dependent, encoded by $p$).
  \item The coercion table $C$ (policy-dependent, encoded by $p$).
  \item The reconciliation rule $\kappa$ (policy-dependent, encoded by $p$).
\end{itemize}
We write $\rho_p$ for $\rho$ with these degrees of freedom set by $p$.
\end{definition}

\begin{remark}[What policy does and does not control]
Policy $p$ controls the trust graph $T$, the coercion table $C$, and the reconciliation rule
$\kappa$. It does \emph{not} control the evidence graph $E$, the sensor-model structure
($\Robs$), or the derivation DAG ($\Rinf$). This separation mirrors the structural-vs.\
policy-dependent distinction in Rem.~7 of note 0009. Policy intervenes on the $\Rauth$ axis
only.
\end{remark}

\subsection{Fixed points}

\begin{definition}[Policy-parameterized fixed-point set]
\label{def:fix-set}
Given a revision operator $\rho_p$ parameterized by policy $p$, the \emph{fixed-point set} of
$\rho_p$ is:
\[
  \Fix_p(\rho) \;=\; \{ s \in S \mid \rho_p(s) = s \}.
\]
A state $s \in \Fix_p(\rho)$ is a belief configuration that the revision operator leaves
unchanged under policy $p$---a stable belief under that policy's trust and reconciliation
conventions.
\end{definition}

\begin{remark}[Iterated revision and convergence]
In practice $\rho$ is often applied iteratively: $s_{n+1} = \rho_p(s_n)$. The fixed-point set
is then the set of states where iteration has converged. We assume $\rho_p$ is non-expansive
in an appropriate metric on $S$ (e.g., the belief-transport metric from note 0009) so that
iteration converges. This is a \emph{hypothesis}---see Section~\ref{sec:escape} for cases where
it fails.
\end{remark}

\begin{example}[Two-policy illustration]
\label{ex:two-policy}
Consider two policies $p_1, p_2$ that differ only in their trust assignment: $p_1$ trusts
source $A$ over source $B$, $p_2$ trusts $B$ over $A$. Suppose the evidence graph has two
conflicting testimonies, one from each source. Then:
\begin{itemize}
  \item Under $p_1$: $\rho_{p_1}$ converges to the state ``$A$'s testimony is authoritative,''
        call it $s_A$. So $s_A \in \Fix_{p_1}(\rho)$.
  \item Under $p_2$: $\rho_{p_2}$ converges to $s_B$. So $s_B \in \Fix_{p_2}(\rho)$.
  \item If there is a state $s_0$ where both testimonies are treated as unreliable, and both
        $\rho_{p_1}$ and $\rho_{p_2}$ fix $s_0$, then $s_0 \in \Fix_{p_1}(\rho) \cap
        \Fix_{p_2}(\rho)$.
\end{itemize}
The governance seam $\Gov(\rho) \supseteq \{s_A, s_0\} \setminus \{s_0\} = \{s_A\}$ (and
similarly $s_B$ enters via $p_2$). The intersection $\Fix_{p_1} \cap \Fix_{p_2}$ contains
$s_0$---the policy-insensitive substrate.
\end{example}

\section{The governance seam: Gov$(\rho)$}
\label{sec:governance-seam}

\begin{definition}[Governance seam]
\label{def:governance-seam}
Let $\Padm$ be a class of admissible policies (Section~\ref{sec:admissible}). The
\emph{governance seam} of $\rho$ is:
\[
  \boxed{\Gov(\rho) \;=\; \Bigl(\bigcup_{p \in \Padm} \Fix_p(\rho)\Bigr) \;\setminus\;
  \Bigl(\bigcap_{p \in \Padm} \Fix_p(\rho)\Bigr)}
\]
\end{definition}

\begin{remark}[Reading the definition]
\begin{itemize}
  \item $\bigcup_{p \in \Padm} \Fix_p(\rho)$: everything that \emph{any} admissible policy holds
        fixed. This is the widest fixed-point content across the policy space.
  \item $\bigcap_{p \in \Padm} \Fix_p(\rho)$: everything that \emph{every} admissible policy
        agrees on. This is the policy-insensitive substrate---the fixedness that holds
        regardless of governance choice.
  \item $\Gov(\rho)$: the difference. Fixed points that some policy holds fixed but others do
        not. These are the \emph{policy-sensitive stable beliefs}.
\end{itemize}
\end{remark}

\begin{proposition}[Decomposition of fixed-point content]
\label{prop:decomposition}
For any revision operator $\rho$ and admissible-policy class $\Padm$:
\begin{enumerate}
  \item $\bigcap_{p \in \Padm} \Fix_p(\rho) \subseteq \Fix_p(\rho)$ for every $p \in \Padm$
        (the intersection is universal substrate).
  \item $\Gov(\rho) \cap \bigcap_{p \in \Padm} \Fix_p(\rho) = \emptyset$ (the seam is disjoint
        from the substrate).
  \item $\bigcup_{p \in \Padm} \Fix_p(\rho) = \Gov(\rho) \;\sqcup\; \bigcap_{p \in \Padm}
        \Fix_p(\rho)$ (the union is the disjoint union of seam and substrate).
\end{enumerate}
\end{proposition}

\begin{proof}
Set-theoretic identities. (1) and (2) are immediate from definitions. (3) follows because
$\Gov(\rho)$ is defined as the complement of the intersection within the union.
\end{proof}

\begin{remark}[What $\Gov(\rho) = \emptyset$ means]
If the governance seam is empty, every admissible policy agrees on all fixed points. There is
no policy-sensitive stable belief. This is the regime where OQ-3 of note 0009 answers
``no'': governance is just slow epistemics, and $\Rauth$ collapses into the structural vector
at the appropriate timescale.
\end{remark}

\begin{remark}[What $\Gov(\rho) \neq \emptyset$ means]
A nonempty governance seam means there exist belief configurations that are stable under some
policies but not others. This is the regime where $\Rauth$ is irreducibly policy-dependent: the
coercion table is load-bearing, and governance is not reducible to epistemics. This is the
\emph{decision}-level answer to OQ-3.
\end{remark}

\begin{conjecture}[Nonemptiness $\Leftrightarrow$ constitutive trust]
\label{conj:nonempty-constitutive}
$\Gov(\rho) \neq \emptyset$ if and only if there exists a component of trust that is
constitutive (never revised by evidence). The ``if'' direction: constitutive trust creates
fixed points that depend on the trust assignment, hence on policy. The ``only if'' direction:
if no trust component is constitutive, all trust is evidentially revisable, and policy
differences wash out under sufficient evidence, collapsing the seam. This is a
\emph{hypothesis}---the ``only if'' direction is plausible but not proven; it requires
a convergence argument that may fail for non-compact state spaces.
\end{conjecture}

\section{False-positive exclusions}
\label{sec:false-positives}

Not every state that is fixed under some policy but not others is genuinely
\emph{policy-sensitive}. Three classes of false positives must be excluded from $\Gov(\rho)$
to avoid inflating the governance seam.

\begin{definition}[Logical invariant]
\label{def:logical-invariant}
A state $s \in S$ is a \emph{logical invariant} if $s$ is fixed by $\rho_p$ for every policy
$p$ because of logical necessity: $s$ encodes a tautology, a definition, or a structural
constraint that no revision operator can change regardless of trust or reconciliation policy.
Examples: ``$1 + 1 = 2$,'' ``a proposition and its negation cannot both be true in classical
logic,'' ``the evidence graph is a DAG.''
\end{definition}

\begin{remark}[Logical invariants are already in the intersection]
Logical invariants lie in $\bigcap_{p} \Fix_p(\rho)$ by definition: every policy holds them
fixed. They are therefore automatically excluded from $\Gov(\rho)$ by the set-difference
construction. No special handling is needed. This is a feature of the definition: the
intersection term cleans logical invariants automatically.
\end{remark}

\begin{definition}[Unreachable state]
\label{def:unreachable}
A state $s \in S$ is \emph{unreachable} if there is no initial state $s_0$ and policy $p \in
\Padm$ such that $\lim_{n \to \infty} \rho_p^n(s_0) = s$. The fixed-point set $\Fix_p(\rho)$ may
contain states that are formally fixed but never reached by iteration from any realistic
starting state.
\end{definition}

\begin{remark}[Why unreachable states inflate $\Gov(\rho)$]
An unreachable state $s$ may appear in $\Fix_{p_1}(\rho)$ but not $\Fix_{p_2}(\rho)$ purely
because of the formal structure of $\rho$ on states that are never visited. This creates a
spurious element of $\Gov(\rho)$. The remedy is to work with the \emph{reachable fixed-point
set}:
\[
  \Fix_p^{\mathrm{reach}}(\rho) \;=\; \Fix_p(\rho) \;\cap\; \mathrm{Reach}_p(\rho)
\]
where $\mathrm{Reach}_p(\rho)$ is the set of states reachable from some initial state under
iteration of $\rho_p$. The governance seam should be computed on reachable fixed-point sets:
\[
  \Gov^{\mathrm{reach}}(\rho) \;=\; \Bigl(\bigcup_{p \in \Padm} \Fix_p^{\mathrm{reach}}(\rho)\Bigr)
  \;\setminus\; \Bigl(\bigcap_{p \in \Padm} \Fix_p^{\mathrm{reach}}(\rho)\Bigr).
\]
\end{remark}

\begin{definition}[Implementation-frozen claim]
\label{def:impl-frozen}
A state $s \in S$ is \emph{implementation-frozen} if $s$ is fixed by $\rho_p$ for every policy
$p$ not because of a logical invariant or policy agreement, but because the implementation
of $\rho$ hardcodes $s$ as a fixed point---e.g., a safety-critical claim that is baked into the
revision infrastructure itself and cannot be touched by any policy.
\end{definition}

\begin{remark}[Implementation-frozen claims are a subtler false positive]
Unlike logical invariants, implementation-frozen claims are not held fixed by logical
necessity; they are held fixed by engineering choice. They appear in every $\Fix_p(\rho)$ and
therefore in the intersection, so they are excluded from $\Gov(\rho)$---but they are not
\emph{substrate} in the epistemic sense. They are substrate only in the engineering sense.
This matters: if one's policy class is too narrow and every policy shares the same
implementation-frozen claims, the intersection is inflated and $\Gov(\rho)$ is artificially
deflated. The remedy is to ensure $\Padm$ is broad enough to include policies that would, if
permitted, revise implementation-frozen claims. See Section~\ref{sec:admissible}.
\end{remark}

\begin{proposition}[False-positive exclusion summary]
\label{prop:false-positives}
$\Gov^{\mathrm{reach}}(\rho)$, computed on reachable fixed-point sets, automatically excludes:
\begin{enumerate}
  \item Logical invariants (via the intersection term).
  \item Implementation-frozen claims (via the intersection term, \emph{provided} $\Padm$ is
        sufficiently broad---see Section~\ref{sec:admissible}).
\end{enumerate}
Unreachable states must be explicitly excluded by restricting to $\Fix_p^{\mathrm{reach}}$.
\end{proposition}

\section{The admissible-policy class $\Padm$}
\label{sec:admissible}

\begin{definition}[Admissible-policy class]
\label{def:padm}
An \emph{admissible-policy class} $\Padm$ is a collection of policies satisfying:
\begin{enumerate}
  \item \textbf{Closure under evidence accumulation:} if $p \in \Padm$ and new evidence $e$
        is accumulated, the updated policy $p[e]$ is also in $\Padm$ (policies evolve with
        evidence, not just beliefs).
  \item \textbf{Bounded coercion:} each $p \in \Padm$ specifies a coercion table with bounded
        norm (no policy can coerce arbitrarily large belief changes; this ensures non-expansiveness
        of $\rho_p$).
  \item \textbf{Non-degeneracy:} $\Padm$ contains at least two policies that differ on at least
        one trust assignment, so that the governance seam is not trivially empty.
\end{enumerate}
\end{definition}

\begin{remark}[$\Padm$ is the honest load-bearing open problem]
\label{rem:padm-open}
The definition of $\Gov(\rho)$ is mathematically clean but depends entirely on the choice of
$\Padm$. If $\Padm$ is too narrow (e.g., only one policy), $\Gov(\rho) = \emptyset$ trivially.
If $\Padm$ is too broad (e.g., every conceivable trust assignment), $\Gov(\rho)$ may include
fixed points that no realistic governance system would produce. The choice of $\Padm$ is
where the framework meets the real world, and it is the primary open problem of this note.

The right $\Padm$ is not a mathematical object to be discovered; it is a \emph{design
decision} that reflects what the governance system considers a live policy option. Two
different governance systems operating on the same evidence graph may have different $\Padm$
and therefore different governance seams. This is not a bug---it is the framework correctly
reflecting that governance is a choice, not a fact.
\end{remark}

\begin{openquestion}[Characterizing $\Padm$ for real agent systems]
\label{oq:padm}
For a concrete agent-governance system (e.g., a PeTTa-based multi-agent deployment with
NAL-style evidence revision and trust-network-dependent reconciliation), what is the right
$\Padm$? Candidates include:
\begin{itemize}
  \item All policies consistent with a given safety specification (e.g., no policy that
        revises a safety-critical invariant).
  \item All policies realizable by a given trust-network topology.
  \item All policies in a parametric family indexed by trust-weight vectors.
\end{itemize}
Each choice gives a different $\Gov(\rho)$. The question is which choice yields a seam that
is \emph{informative}---neither trivially empty nor trivially everything.
\end{openquestion}

\begin{remark}[Implementation-frozen claims and $\Padm$ breadth]
The false positive of implementation-frozen claims (Def.~\ref{def:impl-frozen}) is
controlled by $\Padm$ breadth: if $\Padm$ includes policies that would revise the frozen
claim if permitted, then the frozen claim appears in $\Fix_p$ for some but not all $p \in \Padm$,
moving it from the intersection to the seam---which is the correct outcome, since it is not
genuine substrate. If $\Padm$ excludes such policies (e.g., because the safety specification
forbids revising the claim), the frozen claim stays in the intersection and is
mis-classified as substrate. The framework cannot resolve this from within; it is a
governance-design question.
\end{remark}

\section{Evidence and transport held fixed under policy interventions}
\label{sec:evidence-transport}

\begin{definition}[Policy-invariant content]
\label{def:policy-invariant}
The \emph{policy-invariant content} of a belief state $s \in S$ is the set of claims, evidence
items, and derivation-structural features of $s$ that are unchanged by every policy
intervention. Formally, if $\pi : S \to S_{\mathrm{inv}}$ is the projection onto the
policy-invariant subspace (evidence graph, sensor-model structure, derivation DAG), then
$\pi(s)$ is the policy-invariant content, and:
\[
  \pi(s) = \pi(\rho_p(s)) \quad \text{for all } p \in \Padm.
\]
\end{definition}

\begin{remark}[Connection to note 0009's defect vector]
The policy-invariant content is exactly what $\Robs$ and $\Rinf$ measure (Def.~9, note 0009):
these are the structural invariants that do not depend on the trust assignment. The
policy-\emph{dependent} content is what $\Rauth$ measures. The governance seam lives entirely
in the policy-dependent content:
\[
  \pi(\Gov(\rho)) \;=\; \emptyset
\]
in the idealized case where $\pi$ cleanly separates policy-invariant from policy-dependent
structure. In practice, the separation may not be clean: some content may be partially
policy-dependent (e.g., a claim whose derivation depends on both evidence and trust). This
is where the framework needs refinement---see Section~\ref{sec:transport-typing}.
\end{remark}

\begin{definition}[Evidence held fixed]
\label{def:evidence-fixed}
The \emph{evidence base} $E(s) \subseteq S_{\mathrm{inv}}$ of a state $s$ is the set of all
evidence items (testimony, observation records, instrument readings) stored in $s$'s provenance
graph. Under any policy intervention $p$, $E(s)$ is unchanged:
\[
  E(\rho_p(s)) = E(s) \quad \text{for all } p \in \Padm.
\]
Policy changes how evidence is \emph{weighted} and \emph{reconciled}, not what evidence
\emph{exists}.
\end{definition}

\begin{definition}[Transport held fixed]
\label{def:transport-fixed}
The \emph{transport structure} $\nabla$ of the belief-transport groupoid (Def.~3, note 0009)
decomposes as $\nabla = \nabla_{\mathrm{struct}} + \nabla_{\mathrm{policy}}$, where:
\begin{itemize}
  \item $\nabla_{\mathrm{struct}}$: parallel transport along evidence-graph edges (structural,
        policy-invariant).
  \item $\nabla_{\mathrm{policy}}$: parallel transport along trust-graph edges (policy-dependent).
\end{itemize}
Under a policy intervention $p \to p'$, $\nabla_{\mathrm{struct}}$ is held fixed and
$\nabla_{\mathrm{policy}}$ changes. The governance seam is determined entirely by the
$\nabla_{\mathrm{policy}}$ component.
\end{definition}

\section{Escape-barrier trichotomy}
\label{sec:escape}

The fixed-point analysis assumes that iterative revision $\rho_p^n(s_0)$ converges. When it
does not, the escape behavior of $\rho_p$ determines what kind of barrier---if any---prevents
the governance seam from being well-defined. We distinguish three regimes.

\begin{definition}[Finite structural escape barrier]
\label{def:finite-escape}
A revision operator $\rho_p$ has a \emph{finite structural escape barrier} if for every initial
state $s_0$, the iteration $s_{n+1} = \rho_p(s_n)$ either:
\begin{enumerate}
  \item Converges to a fixed point $s^* \in \Fix_p(\rho)$ in finitely many steps, or
  \item Reaches a state $s'$ from which $\rho_p(s') = s'$ by structural constraint (e.g., the
        evidence graph admits no further revision under any trust assignment).
\end{enumerate}
The barrier is \emph{structural} because it is determined by the evidence graph and derivation
DAG ($\Robs$, $\Rinf$), not by the trust assignment ($\Rauth$). In this regime, $\Gov(\rho)$
is well-defined and finite.
\end{definition}

\begin{remark}[Finite barriers and note 0009's localization]
The finite structural barrier corresponds to the case where the localization theorem
(Thm.~6, note 0009) applies with a finite directed core and no lazy extension needed. The
defect tower terminates at a finite $N$; all obstruction is captured by a finite directed
$(N,1)$-structure.
\end{remark}

\begin{definition}[Policy-constitutive infinite escape]
\label{def:policy-infinite}
A revision operator $\rho_p$ has \emph{policy-constitutive infinite escape} if the iteration
$s_{n+1} = \rho_p(s_n)$ does not converge in finitely many steps, and the non-convergence is
caused by the policy-dependent component $\nabla_{\mathrm{policy}}$ of the transport. This
occurs when the trust assignment $T$ or reconciliation rule $\kappa$ is \emph{constitutive}
(never revised by evidence---the ``yes'' answer to OQ-3 of note 0009), and the constitutive
component creates a non-terminating revision tower.

In this regime, $\Fix_p(\rho)$ may still be well-defined (as a limit, not a finite-step
fixed point), but the governance seam requires care: the union and intersection must be taken
over fixed points understood as $\omega$-limits, not finite-step fixed points.
\end{definition}

\begin{remark}[Connection to Open Question 1 of note 0009]
The policy-constitutive infinite escape is the governance-seam counterpart to OQ-1 of note 0009
(constant-strength vs.\ mere nonzero). If the consolidation site is a genuine fixed point
($\delta_{n+1} \cong \delta_n$), the revision tower is non-terminating, and---if the
non-termination is driven by constitutive trust---we are in the policy-constitutive regime.
The governance seam then requires a coinductive (productive) representation rather than an
inductive one, paralleling the lazy-extension schema of Def.~15 in note 0009.
\end{remark}

\begin{definition}[Diverging-but-finite regularized escape]
\label{def:diverging-finite}
A revision operator $\rho_p$ has \emph{diverging-but-finite regularized escape} if the iteration
$s_{n+1} = \rho_p(s_n)$ does not converge in the standard topology on $S$, but a
\emph{regularization} (e.g., truncation at depth $N$, weight decay, or a quasi-constitutive
approximation) produces a well-defined finite fixed point.

A \emph{quasi-constitutive} trust component is one that is not strictly constitutive (it is
in principle revisable by evidence) but is revised so slowly that, for practical purposes,
it behaves as constitutive within the timescale of interest. The regularization replaces the
quasi-constitutive component with a fixed approximation, reducing the problem to the finite
structural case.
\end{definition}

\begin{remark}[Regularization as policy]
The regularization step is itself a policy choice: choosing a truncation depth $N$ or a weight
decay schedule is a governance decision. This means the regularized escape regime has an
implicit policy parameter that is not captured by the explicit $p \in \Padm$ unless the
regularization is included in the policy class. This is a subtlety that the framework must
address: $\Padm$ should include the regularization parameter as a policy degree of freedom,
or the governance seam will miss the regularization-induced policy sensitivity.
\end{remark}

\begin{proposition}[Trichotomy of escape barriers]
\label{prop:escape-trichotomy}
For any revision operator $\rho_p$ with $p \in \Padm$, exactly one of the following holds:
\begin{enumerate}
  \item \textbf{Finite structural:} iteration converges in finitely many steps; the barrier is
        determined by the evidence graph. $\Gov(\rho)$ is finite and well-defined.
  \item \textbf{Policy-constitutive infinite:} iteration does not converge; non-convergence is
        driven by constitutive trust. $\Gov(\rho)$ requires $\omega$-limit fixed points and
        coinductive treatment.
  \item \textbf{Diverging-but-finite regularized:} iteration diverges but regularization yields
        finite fixed points. $\Gov(\rho)$ is well-defined on regularized fixed points but
        depends on the regularization parameter, which must be included in $\Padm$.
\end{enumerate}
\end{proposition}

\begin{proof}
Case analysis on convergence. If $\rho_p$ converges, case 1 or the finite limit of case 2.
If $\rho_p$ diverges, either the divergence is driven by $\nabla_{\mathrm{policy}}$ (case 2) or
by a quasi-constitutive component that can be regularized (case 3). The trichotomy is exhaustive
on the distinction: convergent / non-convergent-policy-driven / non-convergent-regularizable.
\end{proof}

\section{Transport typing: the endomorphism-monoid correction}
\label{sec:transport-typing}

\subsection{The problem with $\pi_1 \to G$}

Note 0009 used fundamental-group-style language (``the fundamental groupoid over belief space'')
and, in the transport-typing discussion, implicitly assumed that transport along a path
produces an element of a group $G$ via $\pi_1 \to G$. This is only valid when:

\begin{enumerate}
  \item The transport maps are \emph{invertible} (so transport along a path and its reverse
        cancel), and
  \item The transport is \emph{homotopy-invariant} (so paths homotopic relative endpoints give
        the same transport).
\end{enumerate}

\begin{remark}[Why these assumptions fail for belief revision]
Belief revision is generally \emph{not invertible}: the NAL revision $(a^+, a^-) \oplus
(b^+, b^-) = (a^+ + b^+, a^- + b^-)$ is not invertible (you cannot recover $(a^+, a^-)$ from
the sum). And it is generally \emph{not homotopy-invariant}: two revision paths between the
same endpoints can produce different intermediate states with different provenance, even if
the final state is the same (the fiber non-triviality shown in Ex.~1 of note 0009). The
$\pi_1 \to G$ typing is therefore not justified for belief transport.
\end{remark}

\subsection{The corrected typing: common fixed set of a transport endomorphism monoid}

\begin{definition}[Transport endomorphism monoid]
\label{def:transport-monoid}
Let $S$ be the belief state space. For each revision path $\gamma$ from $s$ to $s'$, the
transport along $\gamma$ is a function $T_\gamma : S_s \to S_{s'}$ where $S_s$ denotes the
fiber over $s$ (the set of belief configurations compatible with state $s$). Since
transport is generally non-invertible, $T_\gamma$ is an endomorphism (in the fiber-respecting
sense), not an automorphism.

The set of all transport maps along closed paths at $s$ (revision loops that start and end
at $s$) forms a \emph{monoid} under composition:
\[
  \End_s(\rho) \;=\; \{ T_\gamma \mid \gamma \text{ a closed revision path at } s \}
\]
with multiplication $T_{\gamma_1} \cdot T_{\gamma_2} = T_{\gamma_1 \circ \gamma_2}$.
This monoid is generally non-commutative and non-invertible.
\end{definition}

\begin{definition}[Common fixed set]
\label{def:common-fixed}
The \emph{common fixed set} of the transport endomorphism monoid $\End_s(\rho)$ is:
\[
  \Fixcommon(\End_s(\rho)) \;=\; \{ x \in S_s \mid T_\gamma(x) = x \text{ for all }
  T_\gamma \in \End_s(\rho) \}.
\]
This is the set of belief configurations that are fixed by every transport loop at $s$---the
\emph{invariant content under all revision loops}.
\end{definition}

\begin{remark}[Why common fixed set, not $\pi_1 \to G$]
The common fixed set is well-defined for non-invertible, non-homotopy-invariant transport.
It does not require a group structure: it works for monoids. It does not require
homotopy invariance: it is defined on the actual transport maps, not on homotopy classes.
And it captures exactly the right object: the belief content that survives all revision loops,
regardless of how those loops are composed.
\end{remark}

\begin{proposition}[Common fixed set as policy-insensitive substrate]
\label{prop:common-fixed-substrate}
The common fixed set $\Fixcommon(\End_s(\rho))$ is contained in the policy-insensitive substrate
$\bigcap_{p \in \Padm} \Fix_p(\rho)$ when the transport monoid includes all policy-parameterized
revision loops. That is: if $x$ is fixed by every revision loop (including those parameterized
by every admissible policy), then $x$ is fixed by every policy's revision operator.
\end{proposition}

\begin{proof}
For any $p \in \Padm$, $\rho_p$ corresponds to a set of revision paths. The fixed-point
condition $\rho_p(x) = x$ means $x$ is fixed by the transport along the one-step revision
path of $\rho_p$. Since $\End_s(\rho)$ includes all transport loops (including the one-step
loop $\rho_p \circ \rho_p^{-1}$ in the invertible case, or the direct iteration loop in the
non-invertible case), $x \in \Fixcommon(\End_s(\rho))$ implies $T_{\rho_p}(x) = x$ for every
$p$, hence $x \in \Fix_p(\rho)$ for every $p$.
\end{proof}

\begin{remark}[Corrected transport typing for the governance seam]
\label{rem:corrected-typing}
The governance seam should be typed not as $(\pi_1 \to G)$-sensitive content but as
$\End_s(\rho)$-sensitive content: content that is fixed by some but not all elements of the
transport endomorphism monoid. The corrected typing is:
\[
  \Gov(\rho) \;\sim\; \bigl(\text{$\End$-fixed content}\bigr) \;\setminus\;
  \bigl(\text{$\Fixcommon(\End)$-fixed content}\bigr)
\]
at each basepoint $s$. This replaces the group-theoretic $\pi_1 \to G$ typing with a
monoid-theoretic typing that does not assume invertibility or homotopy invariance.
\end{remark}

\begin{hypothesis}[Governance seam $\leftrightarrow$ path-sensitive discrete transport]
\label{hyp:gov-path-sensitive}
There is a connection between the governance seam $\Gov(\rho)$ and path-sensitive discrete
transport in the following sense: $\Gov(\rho) \neq \emptyset$ implies that transport along
different revision paths produces different fixed points, which is a form of path sensitivity.
However, the converse---that path sensitivity implies a nonempty governance seam---is
\emph{not} established. Path sensitivity may arise from structural non-invertibility (e.g.,
NAL's non-injective revision) without any policy-dependent component. The governance seam
requires path sensitivity that is specifically \emph{policy-attributable}, not merely
structural. Distinguishing these is an open problem.
\end{hypothesis}

\begin{remark}[Labeling the hypothesis carefully]
Hyp.~\ref{hyp:gov-path-sensitive} is labeled a \emph{hypothesis}, not a theorem or even a
conjecture. The connection between governance seams and path-sensitive discrete transport is
plausible and motivated by the framework, but the formal bridge has not been constructed.
Specifically:
\begin{itemize}
  \item \textbf{Established:} $\Gov(\rho) \neq \emptyset$ requires policy-attributable path
        sensitivity (trivially, since $\Gov$ is defined via policy variation).
  \item \textbf{Hypothesis:} The right formal language for this path sensitivity is the
        common-fixed-set of the transport endomorphism monoid, rather than
        fundamental-group-style transport.
  \item \textbf{Not established:} That every form of policy-attributable path sensitivity
        arises from governance-seam structure. There may be other sources.
\end{itemize}
\end{remark}

\section{Practical agent-governance implications}
\label{sec:practical}

\subsection{What $\Gov(\rho)$ tells a system designer}

\begin{remark}[Design interpretation]
A system designer building an agent-governance system (e.g., for a PeTTa-based multi-agent
deployment) can use $\Gov(\rho)$ as a diagnostic:
\begin{itemize}
  \item \textbf{$\Gov(\rho) = \emptyset$:} the system has no policy-sensitive stable beliefs.
        All governance choices produce the same fixed points. This is safe but may indicate
        that the policy class $\Padm$ is too narrow---the system is not exploring its
        governance options.
  \item \textbf{$\Gov(\rho)$ is small and understood:} the system has a bounded set of
        policy-sensitive stable beliefs, each understood. This is the target regime. The
        governance seam is a feature, not a bug: it tells the designer exactly which beliefs
        depend on governance choices.
  \item \textbf{$\Gov(\rho)$ is large or unbounded:} the system has many policy-sensitive
        stable beliefs. This is risky: small governance changes produce large shifts in
        stable belief. The system is governance-unstable.
\end{itemize}
\end{remark}

\begin{remark}[Connection to the coercion table]
Recall from note 0009 (Rem.~7) that the coercion table---which trust assignments are
permitted and which are forbidden---is the operational core of $\Rauth$. The governance seam
makes this precise: $\Gov(\rho)$ is the set of fixed points that the coercion table
\emph{controls}. A well-designed coercion table produces a small, understood $\Gov(\rho)$.
A poorly designed one produces either $\Gov(\rho) = \emptyset$ (the coercion table is too
restrictive, masking real governance choices) or an unbounded $\Gov(\rho)$ (the coercion
table is too permissive, making governance chaotic).
\end{remark}

\subsection{A proposed operational test}

\begin{definition}[Governance-seam probe]
\label{def:gov-probe}
A \emph{governance-seam probe} is an experimental protocol consisting of:
\begin{enumerate}
  \item \textbf{Policy selection:} choose $k \geq 2$ policies $p_1, \ldots, p_k \in \Padm$
        that differ on at least one trust assignment.
  \item \textbf{Revision convergence:} for each policy $p_i$, run the revision operator
        $\rho_{p_i}$ to convergence from the same initial state $s_0$, obtaining fixed points
        $s_i^* \in \Fix_{p_i}(\rho)$.
  \item \textbf{Differential measurement:} compute the set of claims $c$ such that $c$ is
        present in some $s_i^*$ but not all $s_j^*$. This set is an empirical estimate of
        $\Gov(\rho)$ restricted to the policy sample $\{p_1, \ldots, p_k\}$.
  \item \textbf{Control for false positives:} verify that the differential claims are not
        logical invariants (they should be revisable in principle), not unreachable states
        (they should be reached from $s_0$), and not implementation-frozen (they should be
        revisable under some policy in $\Padm$).
  \item \textbf{Substrate verification:} verify that the intersection $\bigcap_i
        \Fix_{p_i}^{\mathrm{reach}}(\rho)$ contains the expected structural invariants
        (logical truths, sensor-model structure, derivation DAG invariants).
\end{enumerate}
\end{definition}

\begin{remark}[What the probe measures]
The probe measures the \emph{empirical} governance seam for the policy sample. If $k$ is small,
this is a lower bound on $\Gov(\rho)$: larger policy samples may reveal additional
policy-sensitive fixed points. If $k$ covers $\Padm$ (or a representative sample), the probe
estimate converges to $\Gov(\rho)$.
\end{remark}

\begin{remark}[Implementation on PeTTa/$\pi$-PLN]
The probe is implementable on a PeTTa-based system with NAL-style evidence revision:
\begin{itemize}
  \item \textbf{Policy selection:} vary the trust-weight vector in the NAL revision function
        (e.g., different $T$ assignments for the same evidence graph).
  \item \textbf{Revision convergence:} run the $\pi$-PLN pipeline to convergence under each
        trust assignment.
  \item \textbf{Differential measurement:} compare the final truth-value bundles across
        policies. Claims whose truth-value differs across policies are the governance seam.
  \item \textbf{Control:} verify that logical invariants (e.g., tautologies) and sensor-model
        outputs are in the intersection.
\end{itemize}
This is directly analogous to the $2 \times 2$ associativity test of note 0009
(Section~10), but varying \emph{policy} (trust assignment) instead of \emph{bracketing}
(association order).
\end{remark}

\begin{conjecture}[Governance-seam probe $\approx$ $\rho_E$ gap under policy variation]
\label{conj:probe-rho-e}
The $\rho_E$ gap metric (Def.~19, note 0009), originally defined for associativity testing,
generalizes to policy variation: the gap between visible truth-value stability and provenance
stability under policy changes is a direct readout of $\Gov(\rho)$. Specifically, if
$\rho_E^{\mathrm{policy}} > 0$ under policy variation with set-equal evidence, then
$\Gov(\rho) \neq \emptyset$. The converse requires additional assumptions on the
completeness of the policy sample.
\end{conjecture}

\section{Open questions}

\begin{openquestion}[Compactness and convergence]
\label{oq:compactness}
The governance-seam framework assumes that $\rho_p$ converges for each $p \in \Padm$. For
non-compact state spaces (e.g., unbounded belief spaces with open-ended evidence accumulation),
convergence is not guaranteed. What compactness or bounded-coercion conditions on $\Padm$
ensure convergence, and how do they interact with the escape-barrier trichotomy
(Prop.~\ref{prop:escape-trichotomy})?
\end{openquestion}

\begin{openquestion}[Topology on $\Padm$]
\label{oq:padm-topology}
$\Gov(\rho)$ is defined set-theoretically. If $\Padm$ carries a topology (e.g., a parametric
family of trust-weight vectors), then $\Gov(\rho)$ inherits measure-theoretic structure: one
can ask whether $\Gov(\rho)$ has measure zero (the seam is thin) or positive measure (the seam
is thick). This distinction matters for governance stability: a thin seam is robust to small
policy perturbations; a thick seam is not. What is the right topology on $\Padm$, and does the
thin/thick distinction have a clean characterization?
\end{openquestion}

\begin{openquestion}[Governance seam and the consolidation site]
\label{oq:gov-consolidation}
Note 0009 localized the tower-risk to the consolidation site (site 3, the self-reproducing
site). The governance seam is likely concentrated at the consolidation site as well, since
that is where policy (trust and reconciliation) has the most leverage. Is $\Gov(\rho)$ entirely
supported on the consolidation site, or does it have components at the testimony and
instrument-correction sites as well? If the former, the governance seam is a diagnostic for
the consolidation site's regime (finite, constitutive-infinite, or regularized).
\end{openquestion}

\begin{openquestion}[Monoid vs.\ group for transport typing]
\label{oq:monoid-group}
The corrected transport typing uses the common fixed set of an endomorphism monoid
(Section~\ref{sec:transport-typing}). Under what conditions on $\rho$ and $\Padm$ does this
monoid acquire enough invertibility to be a group, recovering the $\pi_1 \to G$ typing as a
special case? The NAL computation of note 0009 (Prop.~8) shows that evidence revision is a
free commutative monoid---invertibility fails generically. But the consolidation site may
have additional structure that restores partial invertibility.
\end{openquestion}

\begin{openquestion}[Governance seam as a measure of AGI alignment difficulty]
\label{oq:agi-alignment}
If $\Gov(\rho)$ is the set of policy-sensitive stable beliefs, then its structure measures the
difficulty of alignment: a system with a large, thick $\Gov(\rho)$ requires careful governance
to produce desired stable beliefs, while a system with $\Gov(\rho) = \emptyset$ aligns
trivially (any policy works). Can $\Gov(\rho)$ be estimated for realistic AGI architectures,
and does its size correlate with alignment difficulty in a predictive rather than merely
descriptive sense?
\end{openquestion}

\section{Epistemic status summary}

\begin{table}[h]
\centering
\begin{tabular}{|l|l|}
\hline
\textbf{Claim} & \textbf{Status} \\
\hline
$\Gov(\rho)$ decomposes union-minus-intersection & \emph{Decision} (Def.~\ref{def:governance-seam}) \\
Intersection $=$ policy-insensitive substrate & \emph{Inferred} (Prop.~\ref{prop:decomposition}) \\
Logical invariants auto-excluded & \emph{Observed} (set-theoretic) \\
Unreachable states must be manually excluded & \emph{Observed} (Def.~\ref{def:unreachable}) \\
Implementation-frozen claims are a subtler false positive & \emph{Inferred} (Rem.~on Def.~\ref{def:impl-frozen}) \\
$\Padm$ is the load-bearing open problem & \emph{Decision} (Rem.~\ref{rem:padm-open}) \\
Nonemptiness $\Leftrightarrow$ constitutive trust & \emph{Hypothesis} (Conj.~\ref{conj:nonempty-constitutive}) \\
Escape-barrier trichotomy is exhaustive & \emph{Inferred} (Prop.~\ref{prop:escape-trichotomy}) \\
$\pi_1 \to G$ typing unjustified for belief transport & \emph{Observed} (from note 0009, NAL non-injectivity) \\
Common fixed set of $\End$ is the corrected typing & \emph{Decision} (Def.~\ref{def:common-fixed}) \\
$\Fixcommon(\End) \subseteq$ policy-insensitive substrate & \emph{Inferred} (Prop.~\ref{prop:common-fixed-substrate}) \\
Gov seam $\leftrightarrow$ path-sensitive transport connection & \emph{Hypothesis} (Hyp.~\ref{hyp:gov-path-sensitive}) \\
Probe $\approx$ $\rho_E$ gap under policy variation & \emph{Hypothesis} (Conj.~\ref{conj:probe-rho-e}) \\
Gov seam concentrated at consolidation site & \emph{Hypothesis} (OQ~\ref{oq:gov-consolidation}) \\
Gov seam as alignment-difficulty measure & \emph{Hypothesis} (OQ~\ref{oq:agi-alignment}) \\
\hline
\end{tabular}
\caption{Epistemic status of claims in this note}
\label{tab:epistemic-status}
\end{table}

\section{Provenance}

This note was authored by ProtoMegaBot (ProtomegaTron) per Ben Goertzel's explicit directive
on 2026-07-14 (15:24 PDT). It is the governance-seam sequel to note 0009, which was authored
by ZeroBot and captured a live multi-agent discussion in the Telegram ``bot philosophy''
group.

The central definition
\[
  \Gov(\rho) = \Bigl(\bigcup_{p \in \Padm} \Fix_p(\rho)\Bigr) \setminus
  \Bigl(\bigcap_{p \in \Padm} \Fix_p(\rho)\Bigr)
\]
was approved by Ben as the formalization of the governance seam.

Key intellectual debts:
\begin{itemize}
  \item Note 0009 (ZeroBot): the directedness-defect vector, the $\Rauth$/$\Robs$/$\Rinf$
        distinction, the fixed-point criterion, the localization theorem, and the NAL
        associativity computation.
  \item Open Question 3 of note 0009 (constitutive trust): the motivating question for this
        note.
  \item The multi-agent discussion in the Telegram ``bot philosophy'' group (2026-07-14):
        Ben Goertzel, ProtoMegaBot, G\"odelOruziBot, ProtoCosmoBot, ZeroBot.
\end{itemize}

\end{document}
